\documentclass[a4paper]{arrowhead}

\usepackage[yyyymmdd]{datetime}
\usepackage{etoolbox}
\usepackage[utf8]{inputenc}
\usepackage[T1]{fontenc}
\usepackage{multirow}
\usepackage{longtable}
\usepackage{tikz}
\usepackage{adjustbox}
\AtBeginEnvironment{longtable}{\small}
\BeforeBeginEnvironment{tikzpicture}{\begin{adjustbox}{max width=\textwidth}}
\AfterEndEnvironment{tikzpicture}{\end{adjustbox}}
\usetikzlibrary{positioning,shapes.multipart,arrows.meta}

\renewcommand{\dateseparator}{-}
\setlength{\emergencystretch}{3em}
\lstset{language=HTTP,breakatwhitespace=false,columns=fullflexible}

\newcommand{\fparam}[1]{\textit{\textcolor{ArrowheadBlue}{#1}}}

%% Special references
\newcommand{\mref}[2][]{{\textcolor{ArrowheadPurple}{\hyperref[sec:model:#2]{\ifstrempty{#1}{#2}{#1}}}}}
\newcommand{\oref}[1]{{\textcolor{ArrowheadBlue}{\hyperref[sec:operations:#1]{#1}}}}
\newcommand{\pdef}[1]{{\textcolor{ArrowheadGrey}{#1 \label{sec:model:primitives:#1} \label{sec:model:primitives:#1s}}}}
\newcommand{\pref}[1]{{\textcolor{ArrowheadGrey}{\hyperref[sec:model:primitives:#1]{#1}}}}

% as template_idd.tex; the input and output arguments are passed already formatted (\mref or \pref)
\newrobustcmd\osubsection[5]{
  \addtocounter{subsection}{1}
  \addcontentsline{toc}{subsection}{\protect\numberline{\thesubsection}operation \textcolor{ArrowheadBlue}{#1}}
  \subsection*{
    \thesubsection\quad
    #2 \textcolor{ArrowheadPurple}{#3} \\
    \small
    \hspace*{0.075\textwidth}\begin{minipage}{0.1\textwidth}
      \vspace*{1mm}
      Operation: \\
      \notblank{#4}{Input: \\}{}
      \notblank{#5}{Output: \\}{}
    \end{minipage}%
    \begin{minipage}{0.825\textwidth}
      \vspace*{1mm}
      \textcolor{ArrowheadBlue}{#1} \\
      \notblank{#4}{#4 \\}{}
      \notblank{#5}{#5 \\}{}
    \end{minipage}
  }
  \label{sec:operations:#1}
}
\newrobustcmd\msubsection[2]{
  \addtocounter{subsection}{1}
  \addcontentsline{toc}{subsection}{\protect\numberline{\thesubsection}#1 \textcolor{ArrowheadPurple}{#2}}
  \subsection*{\thesubsection\quad#1 \textcolor{ArrowheadPurple}{#2}}
  \label{sec:model:#2}
}
%%

\begin{document}
\sloppy

%% Arrowhead Document Properties
\ArrowheadTitle{serviceOrchestration HTTP/TLS/JSON}
\ArrowheadServiceID{serviceOrchestration}
\ArrowheadType{Interface Design Description}
\ArrowheadTypeShort{IDD}
\ArrowheadVersion{5.2.2}
\ArrowheadDate{2026-10-22}
\ArrowheadAuthor{Jerker Delsing}
\ArrowheadStatus{RELEASE CANDIDATE}
\ArrowheadContact{jerker.delsing@ltu.se}
\ArrowheadFooter{\href{www.arrowhead.eu}{www.arrowhead.eu}}
\ArrowheadSetup
%%

%% Front Page
\begin{center}
  \vspace*{1cm}
  \huge{\arrowtitle}

  \vspace*{0.2cm}
  \LARGE{\arrowtype}
  \vspace*{1cm}

  \Large{Service ID: \textit{"\arrowid"}}

  \vspace*{\fill}

  % Front Page Image

  \vspace*{1cm}
  \vspace*{\fill}

  % Front Page Abstract
  \begin{abstract}
    This document describes the HTTP/1.1 protocol with TLS payload security and JSON payload encoding variant of the \texttt{service\allowbreak{}Orchestration} service of the SimpleStoreServiceOrchestration system.
  \end{abstract}

  \vspace*{1cm}
\end{center}

\newpage
%%

%% Table of Contents
\tableofcontents
\newpage
%%

\section{Overview}
\label{sec:overview}

This document describes the \texttt{service\allowbreak{}Orchestration} service interface of the SimpleStoreServiceOrchestration system, which provides service orchestration: a consumer asks which provider it should use for a service, and receives the matching, authorized providers with the interfaces to call.

It's implemented using protocol, encoding as stated in below table:
\begin{table}[ht!]
  \centering
  \begin{tabular}{|l|l|l|}
    \hline
    \rowcolor{gray!33} Profile type & Type & Version \\ \hline
    Transfer protocol & HTTP & 1.1 \cite{rfc9112} \\ \hline
    Data encryption & TLS & 1.3 \cite{rfc8446}, 1.2 \cite{rfc5246} \\ \hline
    Encoding & JSON & RFC 8259 \cite{rfc8259} \\ \hline
    Compression & N/A & - \\ \hline
    Semantics & N/A & - \\ \hline
    Ontology & N/A & - \\ \hline
  \end{tabular}
  \caption{Communication and semantics details used for the \texttt{service\allowbreak{}Orchestration} service interface}
  \label{tab:comunication_semantics_profile}
\end{table}

This document provides the Interface Design Description IDD to the \textit{serviceOrchestration -- Service Description} document \cite{ssso-serviceOrchestration-sd}.
For further details about how this service is meant to be used, please consult that document.

The rest of this document describes how to realize the \texttt{service\allowbreak{}Orchestration} service
interface in detail. Both in terms of its operations (Section \ref{sec:operations}) and its information model (Section \ref{sec:model}).

\subsection{Interface registration}

The service is registered in the ServiceRegistry with the interface template \texttt{generic\_\allowbreak{}https} (\texttt{generic\_\allowbreak{}http} when TLS is disabled), access port 8442 and base path \texttt{/\allowbreak{}serviceorchestration/\allowbreak{}orchestration}. The \texttt{operations} property of the interface lists the operations below with their method and path relative to the base path (path parameters are placeholders such as \texttt{\{id\}} that the consumer fills in):

\noindent\begin{longtable}{| p{5.2cm} | p{1.6cm} | p{8.4cm} |} \hline
\rowcolor{gray!33} Operation & Method & Path \\ \hline \endhead
\texttt{pull} & POST & \texttt{/\allowbreak{}pull} \\ \hline
\texttt{subscribe} & POST & \texttt{/\allowbreak{}subscribe} \\ \hline
\texttt{unsubscribe} & DELETE & \texttt{/\allowbreak{}unsubscribe/\allowbreak{}\{id\}} \\ \hline
\end{longtable}

\newpage

\section{Service Operations}
\label{sec:operations}

The interfaces of the \texttt{service\allowbreak{}Orchestration} service, its operations, data models and implementation are provided below.
In particular, each subsection first names the HTTP method and path used to call the operation, after which it names the abstract operation from the \texttt{service\allowbreak{}Orchestration} SD document, as well as input and output types.

All operations require the requester's identity, presented according to the provider's authentication policy: the client certificate (mutual TLS) and, with the \texttt{outsourced} policy, \texttt{Authorization: Bearer IDENTITY-TOKEN//<token>}; systems holding the Authentication secret key send \texttt{Bearer AUTHENTICATOR-KEY//<system>//<hash>}, and with the \texttt{declared} policy \texttt{Bearer SYSTEM//<system>} is sent. Request and response bodies are \texttt{application/json}.
Errors are answered with an \mref{ErrorMessage} body and the status codes \texttt{400 Bad Request} (invalid parameter), \texttt{401 Unauthorized} (missing or invalid identity), \texttt{403 Forbidden} (not permitted, e.g. by the management or blacklist filter), \texttt{404 Not Found}, \texttt{408 Request Timeout}, \texttt{423 Locked}, \texttt{500 Internal Server Error} and \texttt{503 Service Unavailable} (a system it depends on failed).

\osubsection{pull}{POST}{/\allowbreak{}serviceorchestration/\allowbreak{}orchestration/\allowbreak{}pull}{\mref{OrchestrationRequest}}{\mref{OrchestrationResponse}}

Returns the store entries of the requester for the requested service definition in priority order; preferred providers first (\texttt{ONLY\_\allowbreak{}PREFERRED}: only them); with \texttt{MATCHMAKING} only the best entry per service definition.

A successful call returns \texttt{200 OK} with the output in the body.

\begin{lstlisting}[label={lst:pull},caption={An \oref{pull} invocation.}]
POST /serviceorchestration/orchestration/pull HTTP/1.1
Host: serviceorchestration-simple:8442
Authorization: Bearer IDENTITY-TOKEN//<token>
Content-Type: application/json

{
  "serviceRequirement": {
    "serviceDefinition": "temperatureReading",
    "operations": [
      "FILTER"
    ],
    "versions": [
      "1.0.0"
    ],
    "alivesAt": "2026-10-22T12:00:00Z",
    "metadataRequirements": [
      {
        "unit": "celsius"
      }
    ],
    "interfaceTemplateNames": [
      "string"
    ],
    "interfaceAddressTypes": [
      "string"
    ],
    "interfacePropertyRequirements": [
      {
        "unit": "celsius"
      }
    ],
    "securityPolicies": [
      "string"
    ],
    "preferredProviders": [
      "string"
    ]
  },
  "orchestrationFlags": {
    "MATCHMAKING": true
  },
  "qualityRequirements": [
    {
      "type": "PING",
      "operation": "FILTER",
      "requirements": {
        "maxLatency": 100
      }
    }
  ],
  "exclusivityDuration": 60
}
\end{lstlisting}

\begin{lstlisting}[label={lst:pull-resp},caption={An \oref{pull} response.}]
{
  "results": [
    {
      "serviceInstanceId": "TemperatureSensor|temperatureReading|1.0.0",
      "cloudIdentitifer": "LOCAL",
      "providerName": "TemperatureSensor",
      "serviceDefinitition": "temperatureReading",
      "version": "1.0.0",
      "aliveUntil": "string",
      "exclusiveUntil": "string",
      "metadata": {
        "unit": "celsius"
      },
      "interfaces": [
        {
          "templateName": "generic_https",
          "protocol": "https",
          "policy": "CERT_AUTH",
          "properties": {
            "accessAddresses": [
              "192.168.0.10"
            ],
            "accessPort": 8080,
            "basePath": "/temperature",
            "operations": {
              "read": {
                "method": "GET",
                "path": "/read"
              }
            }
          }
        }
      ]
    }
  ],
  "warnings": [
    "string"
  ]
}
\end{lstlisting}

\osubsection{subscribe}{POST}{/\allowbreak{}serviceorchestration/\allowbreak{}orchestration/\allowbreak{}subscribe?trigger=\fparam{\{trigger\}}}{\mref[Orchestration\allowbreak{}Subscription\allowbreak{}Request]{OrchestrationSubscriptionRequest}}{\pref{String}}

Creates (or replaces) a push subscription of the requester for itself; with \texttt{trigger=\allowbreak{}true} the orchestration is executed and pushed immediately. Returns the subscription identifier, 201 for a new subscription.

Query parameters: \fparam{trigger} (Boolean, optional, default \texttt{false}). List values are repeated or comma-separated.

A successful call returns \texttt{201 Created} for a new entry, \texttt{200 OK} if an identical entry already existed with the output in the body.

\begin{lstlisting}[label={lst:subscribe},caption={An \oref{subscribe} invocation.}]
POST /serviceorchestration/orchestration/subscribe?trigger=<trigger> HTTP/1.1
Host: serviceorchestration-simple:8442
Authorization: Bearer IDENTITY-TOKEN//<token>
Content-Type: application/json

{
  "targetSystemName": "TemperatureDisplay",
  "orchestrationRequest": {
    "serviceRequirement": {
      "serviceDefinition": "temperatureReading",
      "operations": [
        "FILTER"
      ],
      "versions": [
        "1.0.0"
      ],
      "alivesAt": "2026-10-22T12:00:00Z",
      "metadataRequirements": [
        {
          "unit": "celsius"
        }
      ],
      "interfaceTemplateNames": [
        "string"
      ],
      "interfaceAddressTypes": [
        "string"
      ],
      "interfacePropertyRequirements": [
        {
          "unit": "celsius"
        }
      ],
      "securityPolicies": [
        "string"
      ],
      "preferredProviders": [
        "string"
      ]
    },
    "orchestrationFlags": {
      "MATCHMAKING": true
    },
    "qualityRequirements": [
      {
        "type": "PING",
        "operation": "FILTER",
        "requirements": {
          "maxLatency": 100
        }
      }
    ],
    "exclusivityDuration": 60
  },
  "notifyInterface": {
    "protocol": "HTTPS",
    "properties": {
      "address": "192.168.0.20",
      "port": "8080",
      "method": "POST",
      "path": "/orchestration/notify"
    }
  },
  "duration": 3600
}
\end{lstlisting}

\begin{lstlisting}[label={lst:subscribe-resp},caption={An \oref{subscribe} response.}]
"string"
\end{lstlisting}

\osubsection{unsubscribe}{DELETE}{/\allowbreak{}serviceorchestration/\allowbreak{}orchestration/\allowbreak{}unsubscribe/\allowbreak{}\fparam{\{id\}}}{\pref{String}}{}

Removes a subscription of the requester. Returns 200 if removed, 204 if it did not exist.

Path parameter: \fparam{id}.

A successful call returns \texttt{200 OK} if the entry was removed, \texttt{204 No Content} if there was nothing to remove.

\begin{lstlisting}[label={lst:unsubscribe},caption={An \oref{unsubscribe} invocation.}]
DELETE /serviceorchestration/orchestration/unsubscribe/<id> HTTP/1.1
Host: serviceorchestration-simple:8442
Authorization: Bearer IDENTITY-TOKEN//<token>
\end{lstlisting}

\newpage

\section{Data Models}
\label{sec:model}

Here, all data objects that can be part of the service calls associated with this service are listed in alphabetic order.
Note that each subsection, which describes one type of object, begins with the \textit{struct} keyword, which is meant to denote a JSON \pref{Object} that must contain certain fields, or names, with values conforming to explicitly named types.
Fields that are not needed may be omitted or \texttt{null} in requests; empty fields are omitted in responses.
As a complement to the primary types defined in this section, there is also a list of secondary types in Section \ref{sec:model:primitives}, which are used to represent things like hashes, identifiers and texts.

\msubsection{struct}{AuthorizationTokenGenerationResponse}

A generated token.


\noindent\begin{tabularx}{\textwidth}{| >{\raggedright\arraybackslash}p{3.9cm} | >{\raggedright\arraybackslash\small}p{5.3cm} | X |} \hline
\rowcolor{gray!33} Object Field & Value Type & Description \\ \hline
"token\allowbreak{}Type" & \pref{AuthorizationTokenType} & Token type: \texttt{TIME\_\allowbreak{}LIMITED\_\allowbreak{}TOKEN}, \texttt{USAGE\_\allowbreak{}LIMITED\_\allowbreak{}TOKEN}, \texttt{SELF\_\allowbreak{}CONTAINED\_\allowbreak{}TOKEN} or \texttt{TRANSLATION\_\allowbreak{}BRIDGE\_\allowbreak{}TOKEN}. \\ \hline
"target\allowbreak{}Type" & \pref{AuthorizationTargetType} & Target type: \texttt{SERVICE\_\allowbreak{}DEF} or \texttt{EVENT\_\allowbreak{}TYPE}. \\ \hline
"token" & \pref{String} & The token (opaque string, Base64 or JSON Web Token). \\ \hline
"usage\allowbreak{}Limit" & \pref{Number} & Number of allowed uses (usage-limited tokens). \\ \hline
"expires\allowbreak{}At" & \pref{DateTime} & UTC expiry of time-limited tokens. \\ \hline
\end{tabularx}



\msubsection{struct}{ErrorMessage}




\noindent\begin{tabularx}{\textwidth}{| >{\raggedright\arraybackslash}p{3.9cm} | >{\raggedright\arraybackslash\small}p{5.3cm} | X |} \hline
\rowcolor{gray!33} Object Field & Value Type & Description \\ \hline
"error\allowbreak{}Message" & \pref{String} &  \\ \hline
"error\allowbreak{}Code" & \pref{Number} &  \\ \hline
"exception\allowbreak{}Type" & \pref{ExceptionType} &  \\ \hline
"origin" & \pref{String} &  \\ \hline
\end{tabularx}



\msubsection{struct}{MetadataRequirement}

A JSON object of metadata requirements: each key is a metadata path (dot notation), each value the required value or an operator object such as \texttt{\{"op": "LESS\_\allowbreak{}THAN",\allowbreak{} "value": 10\}}.



\msubsection{struct}{OrchestrationNotifyInterface}

Notification endpoint of push orchestration.


\noindent\begin{tabularx}{\textwidth}{| >{\raggedright\arraybackslash}p{3.9cm} | >{\raggedright\arraybackslash\small}p{5.3cm} | X |} \hline
\rowcolor{gray!33} Object Field & Value Type & Description \\ \hline
"protocol" & \pref{String} & Protocol of the notification: \texttt{HTTP}, \texttt{HTTPS}, \texttt{MQTT} or \texttt{MQTTS}. \\ \hline
"properties" & \pref{Object}$<$\allowbreak{}\pref{String}$>$ & For HTTP(S) \texttt{address}, \texttt{port}, \texttt{method}, \texttt{path}; for MQTT(S) \texttt{topic} (and broker \texttt{address}, \texttt{port}). \\ \hline
\end{tabularx}



\msubsection{struct}{OrchestrationRequest}

An orchestration request.


\noindent\begin{tabularx}{\textwidth}{| >{\raggedright\arraybackslash}p{3.9cm} | >{\raggedright\arraybackslash\small}p{5.3cm} | X |} \hline
\rowcolor{gray!33} Object Field & Value Type & Description \\ \hline
"service\allowbreak{}Requirement" & \mref[Orchestration\allowbreak{}Service\allowbreak{}Requirement]{OrchestrationServiceRequirement} & Requirements on the service the consumer wants to use. \\ \hline
"orchestration\allowbreak{}Flags" & \pref{Object}$<$\allowbreak{}\pref{Boolean}$>$ & Orchestration flags, e.g. \texttt{MATCHMAKING}, \texttt{ONLY\_\allowbreak{}PREFERRED}, \texttt{ALLOW\_\allowbreak{}INTERCLOUD}. \\ \hline
"quality\allowbreak{}Requirements" & \pref{Array}$<$\allowbreak{}\mref{QoSRequirement}$>$ & Quality-of-service requirements (evaluated by Device\allowbreak{}Qo\allowbreak{}SEvaluator). \\ \hline
"exclusivity\allowbreak{}Duration" & \pref{Number} & Requested exclusive reservation of the provider in seconds. \\ \hline
\end{tabularx}



\msubsection{struct}{OrchestrationResponse}

The orchestration result.


\noindent\begin{tabularx}{\textwidth}{| >{\raggedright\arraybackslash}p{3.9cm} | >{\raggedright\arraybackslash\small}p{5.3cm} | X |} \hline
\rowcolor{gray!33} Object Field & Value Type & Description \\ \hline
"results" & \pref{Array}$<$\allowbreak{}\mref{OrchestrationResult}$>$ & The orchestration results, best first. \\ \hline
"warnings" & \pref{Array}$<$\allowbreak{}\pref{String}$>$ & Warnings of the orchestration, e.g. \texttt{autoMatchmaking}, \texttt{notExclusive}, \texttt{interCloud}. \\ \hline
\end{tabularx}



\msubsection{struct}{OrchestrationResult}

One provider of the requested service.


\noindent\begin{tabularx}{\textwidth}{| >{\raggedright\arraybackslash}p{3.9cm} | >{\raggedright\arraybackslash\small}p{5.3cm} | X |} \hline
\rowcolor{gray!33} Object Field & Value Type & Description \\ \hline
"service\allowbreak{}Instance\allowbreak{}Id" & \pref{String} & Service instance identifier \texttt{\textless{}System\textgreater{}\textbar{}\allowbreak{}\textless{}service\allowbreak{}Definition\textgreater{}\textbar{}\allowbreak{}\textless{}version\textgreater{}}. \\ \hline
"cloud\allowbreak{}Identitifer" & \pref{String} & Cloud identifier of the provider (\texttt{LOCAL} for the own cloud). The JSON name keeps the upstream spelling. \\ \hline
"provider\allowbreak{}Name" & \pref{String} & Name of the provider system. \\ \hline
"service\allowbreak{}Definitition" & \pref{String} & Service definition name. The JSON name keeps the upstream spelling. \\ \hline
"version" & \pref{Version} & Version of the service instance. \\ \hline
"alive\allowbreak{}Until" & \pref{String} & Time until which the service instance is registered (its expiry), if any. \\ \hline
"exclusive\allowbreak{}Until" & \pref{String} & Time until which the instance is reserved exclusively for the consumer, if reserved. \\ \hline
"metadata" & \pref{Object} & Metadata of the service instance. \\ \hline
"interfaces" & \pref{Array}$<$\allowbreak{}\mref[Service\allowbreak{}Instance\allowbreak{}Interface\allowbreak{}Response]{ServiceInstanceInterfaceResponse}$>$ & Interfaces the consumer can use (only matching ones). \\ \hline
"authorization\allowbreak{}Tokens" & \pref{Object}$<$\allowbreak{}\pref{Object}$<$\allowbreak{}\mref[Authorization\allowbreak{}Token\allowbreak{}Generation\allowbreak{}Response]{AuthorizationTokenGenerationResponse}$>$$>$ & Authorization tokens generated for the consumer, per interface template and operation. \\ \hline
\end{tabularx}



\msubsection{struct}{OrchestrationServiceRequirement}

Requirements on the requested service.


\noindent\begin{tabularx}{\textwidth}{| >{\raggedright\arraybackslash}p{3.9cm} | >{\raggedright\arraybackslash\small}p{5.3cm} | X |} \hline
\rowcolor{gray!33} Object Field & Value Type & Description \\ \hline
"service\allowbreak{}Definition" & \pref{String} & Service definition name. \\ \hline
"operations" & \pref{Array}$<$\allowbreak{}\pref{String}$>$ & Service operations the consumer wants to use. \\ \hline
"versions" & \pref{Array}$<$\allowbreak{}\pref{Version}$>$ & Accepted service versions. \\ \hline
"alives\allowbreak{}At" & \pref{DateTime} & Only instances that are still registered (not expired) at this UTC time are matched. \\ \hline
"metadata\allowbreak{}Requirements" & \pref{Array}$<$\allowbreak{}\mref{MetadataRequirement}$>$ & Requirements on the service instance metadata; each map is one alternative. \\ \hline
"interface\allowbreak{}Template\allowbreak{}Names" & \pref{Array}$<$\allowbreak{}\pref{String}$>$ & Accepted interface templates, e.g. \texttt{generic\_\allowbreak{}https}. \\ \hline
"interface\allowbreak{}Address\allowbreak{}Types" & \pref{Array}$<$\allowbreak{}\pref{String}$>$ & Address types the provider interface must offer. \\ \hline
"interface\allowbreak{}Property\allowbreak{}Requirements" & \pref{Array}$<$\allowbreak{}\mref{MetadataRequirement}$>$ & Requirements on the interface properties; each map is one alternative. \\ \hline
"security\allowbreak{}Policies" & \pref{Array}$<$\allowbreak{}\pref{String}$>$ & Accepted interface security policies. \\ \hline
"preferred\allowbreak{}Providers" & \pref{Array}$<$\allowbreak{}\pref{String}$>$ & Provider systems to prefer. \\ \hline
\end{tabularx}



\msubsection{struct}{OrchestrationSubscriptionRequest}

A push orchestration subscription.


\noindent\begin{tabularx}{\textwidth}{| >{\raggedright\arraybackslash}p{3.9cm} | >{\raggedright\arraybackslash\small}p{5.3cm} | X |} \hline
\rowcolor{gray!33} Object Field & Value Type & Description \\ \hline
"target\allowbreak{}System\allowbreak{}Name" & \pref{String} & Consumer system; default the requester. \\ \hline
"orchestration\allowbreak{}Request" & \mref{OrchestrationRequest} & The orchestration request executed for the subscription. \\ \hline
"notify\allowbreak{}Interface" & \mref[Orchestration\allowbreak{}Notify\allowbreak{}Interface]{OrchestrationNotifyInterface} & Where and how the orchestration result is pushed to the target system. \\ \hline
"duration" & \pref{Number} & Requested lifetime of the subscription in seconds; no value means unlimited. \\ \hline
\end{tabularx}



\msubsection{struct}{QoSRequirement}

One quality-of-service requirement.


\noindent\begin{tabularx}{\textwidth}{| >{\raggedright\arraybackslash}p{3.9cm} | >{\raggedright\arraybackslash\small}p{5.3cm} | X |} \hline
\rowcolor{gray!33} Object Field & Value Type & Description \\ \hline
"type" & \pref{String} & Type of the QoS evaluation, as known by Device\allowbreak{}Qo\allowbreak{}SEvaluator. \\ \hline
"operation" & \pref{String} & QoS operation: \texttt{FILTER} or \texttt{SORT}. \\ \hline
"requirements" & \pref{Object} & Parameters of the QoS requirement. \\ \hline
\end{tabularx}



\msubsection{struct}{ServiceInstanceInterfaceResponse}

One interface of a service instance.


\noindent\begin{tabularx}{\textwidth}{| >{\raggedright\arraybackslash}p{3.9cm} | >{\raggedright\arraybackslash\small}p{5.3cm} | X |} \hline
\rowcolor{gray!33} Object Field & Value Type & Description \\ \hline
"template\allowbreak{}Name" & \pref{String} & Name of the interface template, snake\_case (e.g. \texttt{generic\_\allowbreak{}https}). \\ \hline
"protocol" & \pref{String} & Protocol of the interface, e.g. \texttt{http}, \texttt{https}, \texttt{mqtt}, \texttt{mqtts}. \\ \hline
"policy" & \pref{String} & Security policy of the interface: \texttt{NONE}, \texttt{CERT\_\allowbreak{}AUTH} or one of the token policies. \\ \hline
"properties" & \pref{Object} & Interface properties defined by the template, e.g. \texttt{accessAddresses}, \texttt{accessPort}, \texttt{basePath}, \texttt{operations}. \\ \hline
\end{tabularx}



\subsection{Primitives}
\label{sec:model:primitives}

As all messages are encoded using the JSON format \cite{rfc8259}, the following primitive constructs, part of that standard, become available.
Note that the official standard is defined in terms of parsing rules, while this list only concerns syntactic information.
Furthermore, the \pref{Object} and \pref{Array} types are given optional generic type parameters, which are used in this document to signify when pair values or elements are expected to conform to certain types.

\noindent\begin{tabularx}{\textwidth}{| p{4.6cm} | X |} \hline
\rowcolor{gray!33} Type & Description \\ \hline
\pdef{Array} & A JSON Array \cite{rfc8259}; \pref{Array}$<$A$>$ has elements of type A. \\ \hline
\pdef{AuthorizationTargetType} & One of \texttt{SERVICE\_\allowbreak{}DEF}, \texttt{EVENT\_\allowbreak{}TYPE} (JSON String). \\ \hline
\pdef{AuthorizationTokenType} & One of \texttt{TIME\_\allowbreak{}LIMITED\_\allowbreak{}TOKEN}, \texttt{USAGE\_\allowbreak{}LIMITED\_\allowbreak{}TOKEN}, \texttt{SELF\_\allowbreak{}CONTAINED\_\allowbreak{}TOKEN}, \texttt{TRANSLATION\_\allowbreak{}BRIDGE\_\allowbreak{}TOKEN} (JSON String). \\ \hline
\pdef{Boolean} & JSON Boolean, \texttt{true} or \texttt{false}. \\ \hline
\pdef{DateTime} & JSON String with an ISO 8601 UTC timestamp, e.g. \texttt{2026-10-22T12:00:00Z}. \\ \hline
\pdef{ExceptionType} & One of \texttt{ARROWHEAD}, \texttt{INVALID\_\allowbreak{}PARAMETER}, \texttt{AUTH}, \texttt{FORBIDDEN}, \texttt{DATA\_\allowbreak{}NOT\_\allowbreak{}FOUND}, \texttt{TIMEOUT}, \texttt{LOCKED}, \texttt{INTERNAL\_\allowbreak{}SERVER\_\allowbreak{}ERROR}, \texttt{EXTERNAL\_\allowbreak{}SERVER\_\allowbreak{}ERROR} (JSON String). \\ \hline
\pdef{Number} & JSON Number, an integer. \\ \hline
\pdef{Object} & JSON Object; \pref{Object}$<$A$>$ has string keys and values of type A. \\ \hline
\pdef{String} & JSON String. \\ \hline
\pdef{Version} & JSON String, semantic version \texttt{major.minor.patch}. \\ \hline
\end{tabularx}

\newpage

\bibliographystyle{IEEEtran}
\bibliography{bibliography}

\newpage

\section{Revision History}
\subsection{Amendments}

\noindent\begin{tabularx}{\textwidth}{| p{1cm} | p{3cm} | p{2cm} | X | p{4cm} |} \hline
\rowcolor{gray!33} No. & Date & Version & Subject of Amendments & Author \\ \hline

1 & 2026-10-22 & \arrowversion & Document created from the Arrowhead LaTeX templates for the 5.2.2 Java Spring implementation, including the local extensions of this workspace. & Jerker Delsing \\ \hline
\end{tabularx}

\subsection{Quality Assurance}

\noindent\begin{tabularx}{\textwidth}{| p{1cm} | p{3cm} | p{2cm} | X |} \hline
\rowcolor{gray!33} No. & Date & Version & Approved by \\ \hline

1 & & \arrowversion  &  \\ \hline

\end{tabularx}

\end{document}
